\documentclass[11pt,a4paper]{article}
\usepackage[T1]{fontenc}
\usepackage[utf8]{inputenc}
\usepackage{lmodern,microtype}
\usepackage[margin=27mm,headheight=14pt]{geometry}
\usepackage{amsmath,amssymb,amsthm,mathtools,bm}
\usepackage{booktabs,array,longtable,enumitem}
\usepackage{xcolor,xurl}
\usepackage[colorlinks=true,linkcolor=blue!45!black,citecolor=blue!45!black,urlcolor=blue!45!black]{hyperref}
\usepackage{fancyhdr}
\pagestyle{fancy}\fancyhf{}
\fancyhead[L]{\small Phase averaging in Fabius--Rvachev quadrature}
\fancyhead[R]{\small Research article}
\fancyfoot[C]{\thepage}
\setlength{\emergencystretch}{2em}
\setlength{\parskip}{3pt}
\setcounter{tocdepth}{2}
\numberwithin{equation}{section}
\newtheorem{theorem}{Theorem}[section]
\newtheorem{lemma}[theorem]{Lemma}
\newtheorem{proposition}[theorem]{Proposition}
\newtheorem{corollary}[theorem]{Corollary}
\theoremstyle{definition}
\newtheorem{definition}[theorem]{Definition}
\newtheorem{example}[theorem]{Example}
\newtheorem{question}{Research question}
\theoremstyle{remark}
\newtheorem{remark}[theorem]{Remark}
% ed. (2026-09-30): unnumbered environment for editorial notes.
\newtheorem*{ednote}{Editorial note (ProveIt, 2026-09-30)}
\newcommand{\R}{\mathbb R}\newcommand{\Z}{\mathbb Z}\newcommand{\N}{\mathbb N}
\newcommand{\T}{\mathbb T}\newcommand{\C}{\mathbb C}
\newcommand{\dd}{\,\mathrm d}\newcommand{\ii}{\mathrm i}
\newcommand{\sinc}{\operatorname{sinc}}
\newcommand{\supp}{\operatorname{supp}}\newcommand{\sgn}{\operatorname{sgn}}
\newcommand{\vTwo}{v_2}\newcommand{\TV}{\mathrm{TV}}
\newcommand{\Pol}{\mathcal P}\newcommand{\norm}[1]{\left\lVert#1\right\rVert}
\newcommand{\abs}[1]{\left\lvert#1\right\rvert}
\newcommand{\Q}{\mathcal Q}\newcommand{\Def}{\mathcal D}
\newcommand{\snapshot}{\texttt{37e61c1fdec28c6e7ab7ff445043993077b42cbd}}
\hypersetup{pdftitle={Rigidity and Quantitative Limits of Phase Averaging in Fabius--Rvachev Quadrature},pdfauthor={Research article prepared for Vladimir Reshetnikov},pdfsubject={Signed phase filters, exact quadrature, arithmetic rigidity and superconvergence}}

\begin{document}
% ed. (2026-09-30): no page anchor on the title page (it duplicated page.1).
\hypersetup{pageanchor=false}
\begin{titlepage}
\vspace*{12mm}
\noindent{\large\scshape Analysis / Approximation / Harmonic analysis}\par
\vspace{13mm}
{\LARGE\bfseries Rigidity and Quantitative Limits\par
of Phase Averaging in\par
Fabius--Rvachev Quadrature\par}
\vspace{9mm}
{\large Signed-filter classification, a rational--irrational dichotomy,\par
and sharper superconvergence bounds\par}
\vspace{14mm}
\noindent Prepared for Vladimir Reshetnikov\par
\noindent September 28, 2026\par
\vspace{12mm}
\begin{abstract}
We study quadrature obtained by sampling Rvachev's smooth compactly supported
probability density on a translated lattice and then averaging lattice phases.
For a rational mesh parameter $M=a/b$ in lowest terms, we prove that a finite
signed phase measure integrates every polynomial of degree at most $r$, for
\emph{every} overall lattice translation, precisely when it is invariant under
rotation by $1/L$, where
\[
 L=b\,2^{\max\{0,r-v_2(a)\}}.
\]
Consequently, an atomic rule requires at least $L$ distinct phases. Equality
forces the equally weighted regular $L$-gon, even when negative weights are
allowed. For irrational $M$, the only exact phase measure is Haar measure,
so no finite phase rule is uniformly exact even on constants. A finite
annihilating-polynomial argument strengthens these impossibility statements
to explicit positive error bounds for under-budget signed rules.

A second development improves the repository's odd-dilation estimate for the
Fourier product from an inverse first power to an inverse square, and then
to a log-Gaussian envelope uniform in the odd part. This gives
a uniform relative first-harmonic factorization, a complete integer-mesh
phase classification including the odd-mesh boundary case, simple zeros,
and quantitative phase-error bounds. Existing phase-classification results
in the repository are explicitly credited rather than claimed as new.
All conclusions are proved in ordinary mathematics. Exact rational tests,
independent Fourier checks, a formalization plan, and further research
questions accompany the article. The proposed new refinements and filter
theorems have not been independently peer reviewed or checked in Lean.
\end{abstract}
\vfill
\noindent\small Repository baseline: ProveIt, commit\par
\noindent\small\snapshot.
\end{titlepage}
% ed. (2026-09-30): page anchors resume after the title page.
\hypersetup{pageanchor=true}

\tableofcontents
\newpage

\section{Research target, contributions, and source boundaries}

\subsection{The question that phase averaging raises}

Let $u$ denote the Rvachev up-function normalized to have support $[-1,1]$
and integral one. For $M>0$ and a phase $\theta\in\T=\R/\Z$, consider
\begin{equation}
 \Q_{M,\theta}(P)
 =\frac1M\sum_{k\in\Z}
 P\!\left(\frac{k+\theta}{M}\right)
 u\!\left(\frac{k+\theta}{M}\right).
 \label{eq:quadrature}
\end{equation}
The sum is finite. On integer meshes the repository proves exactness through
degree $v_2(M)$ at every phase, and one additional degree at certain
parity-selected phases. Its prose development also contains stronger
phase-zero classifications and multiple-angle estimates
\cite{repoLean,repoComb}. These results naturally raise a different design
question: \emph{can several carefully chosen phases, possibly with signed
weights, raise the degree of exactness more economically than refining the
lattice?}

There are two distinct quantifiers. A single specially chosen phase can
cancel an error at one translation. A fixed phase filter that must work
\emph{for every overall translation} has to cancel the relevant Fourier
coefficients individually. The main results characterize the latter
requirement exactly. The answer is rigid: for a rational mesh all admissible
phase measures have a prescribed rotational symmetry, and a minimum-size
filter is simply a finer lattice in disguise. Moreover, allowing large
positive and negative weights cannot make an under-budget rule arbitrarily
accurate in the specified uniform moment-error norm.

\subsection{Principal conclusions}

Write $\Pol_r$ for the real polynomials of degree at most $r$, including zero.
For a finite signed Borel measure $\nu$ on $\T$ of mass one, define
\begin{equation}
 \Q^\nu_{M,\theta}(P)
 =\int_\T\Q_{M,\theta+\varphi}(P)\dd\nu(\varphi).
 \label{eq:filtered}
\end{equation}
The following are the main statements, with all proofs supplied below.

\begin{enumerate}[leftmargin=*,label=\textbf{\arabic*.},itemsep=4pt]
\item \textbf{Exact signed-filter classification.}
If $M=a/b\in\Q_{>0}$ is reduced and
$L=b2^{\max(0,r-v_2(a))}$, then \eqref{eq:filtered} is exact on
$\Pol_r$ for every $\theta$ if and only if $\nu$ is invariant under
$\varphi\mapsto\varphi+1/L$ (Theorem~\ref{thm:classification}).
If $M$ is irrational, exactness on constants for every $\theta$ already forces
$\nu$ to be Haar measure (Theorem~\ref{thm:irrational}).
\item \textbf{Optimal cost and equality cases.}
Every exact atomic filter has a number of support points divisible by $L$.
The minimum is $L$, attained only by a rotated uniform $L$-gon.
Negative weights do not improve this minimum
(Theorem~\ref{thm:atomic}). For an integer mesh, each additional
translation-uniform degree requires doubling the minimum phase count.
\item \textbf{A quantitative obstruction below the minimum.}
For $K<L$, an explicit constant $c_{M,r,K}>0$ bounds below the
maximum uniform error on the monomials $1,x,\ldots,x^r$ of \emph{every}
$K$-phase signed rule (Theorem~\ref{thm:quant-obstruction}).
There is an analogous constant-error obstruction for every finite $K$ on
an irrational mesh. Neither result imposes a bound on the weights.
\item \textbf{A sharpened phase theorem.}
For $M=2^dq$, $q$ odd, the first defect at degree $p=d+1$ has a
nonvanishing amplitude times its parity-selected sine or cosine, times
a positive continuous function lying between $1-\eta_p$ and $1+\eta_p$,
where
\[
 \eta_p=(1-2^{-p-1})\zeta(p+1)-1\le\pi^2/8-1<0.234.
\]
This gives explicit root conditioning and a unified treatment of $p=1$
as well as higher degrees (Theorem~\ref{thm:factorization}).
\end{enumerate}

\subsection{What is inherited, and what is proposed here}

The classical Fourier product, its integer zero multiplicities, and the
probability model are established facts; Arias de Reyna's paper is a primary
reference \cite{arias}. Fourier-zero criteria for polynomial reproduction
belong to the Strang--Fix tradition; see, for example, the research treatment
in \cite{zakharov}. We rederive the particular periodization criterion needed
here, rather than assuming that a general reproduction theorem applies
without checking its normalization.

The pinned repository's module
\path{RvachevSuperconvergentSynthesis.lean} explicitly does not assert that its
selected phases are the only ones. This is a limitation of that module's
stated surface, \emph{not} evidence that the entire repository lacks a
phase-classification proof. Indeed, the canonical comb manuscript contains
\path{thm:phase-zero-set}, including general-odd-part corrected Fabius rules
and dyadic Rvachev combs, together with the estimate
$|H(qn/2)|\le |H(q/2)|/n$ \cite{repoComb}. We therefore do not claim to
solve the already-settled dyadic phase conjecture again.

The contribution proposed here is the measure-level classification and its
sharp and quantitative cost consequences, together with the inverse-square
and uniform log-Gaussian improvements, root estimates, and their synthesis. The stronger
comparison uses an additional exactly controlled sinc factor. The proof
works down to $p=1$, where the weaker comparison alone does not give the
required summable relative sine bound. These are mathematically proved
statements in this article, but their worldwide bibliographic novelty is not
certified by the targeted literature and repository review. No claim of a
solution to a major named open problem is made.

\section{The kernel and the exact Fourier criterion}

\subsection{Normalization and elementary analytic facts}

Our Fourier transform convention is
\begin{equation}
 \widehat f(t)=\int_\R f(x)e^{-2\pi\ii tx}\dd x,
 \qquad
 \sinc(t)=\begin{cases}\sin(\pi t)/(\pi t),&t\ne0,\\1,&t=0.\end{cases}
 \label{eq:fourier-convention}
\end{equation}
Let $U_0,U_1,\ldots$ be independent random variables uniform on
$[-1/2,1/2]$, and set
\begin{equation}
 X=\sum_{j\ge0}2^{-j}U_j,
 \qquad
 H(t)=\prod_{j\ge0}\sinc(t/2^j).
 \label{eq:product}
\end{equation}
The series converges absolutely for every outcome and takes values in
$[-1,1]$. Its distribution has characteristic function $H$. Its density is
$u$. This agrees with the repository normalization and with the standard
up-function \cite{arias,repoComb}.

\begin{lemma}[Smoothness, product convergence, and zeros]
\label{lem:kernel}
The distribution of $X$ has an even nonnegative density
$u\in C_c^\infty(\R)$, with support contained in $[-1,1]$ and integral one.
The product $H$ extends to an entire function and converges normally on
compact subsets of $\C$. Its real zeros are exactly the nonzero integers.
At $m\in\Z\setminus\{0\}$ the order of the zero is
\begin{equation}
 \operatorname{ord}_{t=m}H(t)=1+v_2(|m|).
 \label{eq:zero-order}
\end{equation}
In particular $H(q/2)\ne0$ for every odd integer $q$.
\end{lemma}

\begin{proof}
On each fixed compact set in $\C$,
$\sinc(z/2^j)=1+O(4^{-j})$ uniformly. Thus the product is normally
convergent. At a point where none of the factors vanishes, the summable
perturbation from one implies that the tail product is nonzero. A factor
vanishes precisely when $z$ is a nonzero integer multiple of $2^j$.
At a nonzero integer $m$, exactly the factors indexed by
$0\le j\le v_2(|m|)$ vanish, and every such zero is simple. This proves
\eqref{eq:zero-order}, as well as the absence of any other zeros.

For real $t\ne0$ and every positive integer $N$, retaining the first $N$
factors and bounding the rest by one gives
\begin{equation}
 |H(t)|\le
 \frac{2^{N(N-1)/2}}{(\pi|t|)^N}.
 \label{eq:rough-decay}
\end{equation}
Consequently $t^rH(t)$ is integrable for every $r\ge0$. Fourier inversion
of the probability distribution therefore gives a smooth density, with
all derivatives obtained by differentiating the inverse integral. The
support statement follows from the almost-sure bound on $X$; symmetry
follows from the symmetric summands. Nonnegativity and normalization
follow because this is a probability density. All its derivatives vanish
at the endpoints of its support.
\end{proof}

The value $u(0)=1$ can also be seen without a special-function identity:
write $X=U_0+Y$, where $Y$ is supported on $[-1/2,1/2]$. Convolving the
unit-height density of $U_0$ with the distribution of $Y$ gives value one
at zero; the boundary points of $Y$ have probability zero. This observation
also fixes the central sample in the exact computations below.

\subsection{Periodized moment defects}

For $r\ge0$ define
\begin{equation}
 \mu_r=\int_\R x^ru(x)\dd x,
 \qquad
 \Def_{M,r}(\theta)=\Q_{M,\theta}(x^r)-\mu_r.
 \label{eq:defect}
\end{equation}
The finite sum in \eqref{eq:quadrature} is a smooth periodic function of
$\theta$: terms entering or leaving the support have every derivative zero.
The following formula keeps all scaling and sign conventions explicit.

\begin{lemma}[Alias formula]
\label{lem:alias}
For every $M>0$ and $r\ge0$,
\begin{equation}
 \Def_{M,r}(\theta)
 =\left(\frac{\ii}{2\pi}\right)^r
  \sum_{n\in\Z\setminus\{0\}}H^{(r)}(Mn)e^{2\pi\ii n\theta}.
 \label{eq:alias}
\end{equation}
The series and all of its termwise derivatives in $\theta$ converge
absolutely and uniformly.
\end{lemma}

\begin{proof}
Put $f_r(x)=x^ru(x)$. The $n$th Fourier coefficient of the periodic
sum in \eqref{eq:quadrature} is, by substituting $x=(k+\theta)/M$,
\begin{align*}
 \int_0^1\frac1M\sum_k f_r((k+\theta)/M)e^{-2\pi\ii n\theta}\dd\theta
 &=\sum_k\int_{k/M}^{(k+1)/M}f_r(x)e^{-2\pi\ii nMx}\dd x\\
 &=\widehat f_r(Mn).
\end{align*}
There is no extra power of $M$ in this coefficient. Differentiating
$H$ under its integral gives
$\widehat f_r(t)=(\ii/(2\pi))^rH^{(r)}(t)$.
Since $f_r$ is smooth and compactly supported, repeated integration by
parts shows that $\widehat f_r$ decreases faster than any fixed inverse
power. This proves all convergence assertions and identifies the Fourier
series. Its zero coefficient is $\mu_r$, so subtracting it yields
\eqref{eq:alias}.
\end{proof}

\begin{corollary}[Unfiltered, translation-uniform exactness]
\label{cor:unfiltered}
An unfiltered rule at mesh $M>0$ integrates $\Pol_r$ for every phase if and
only if $M$ is a positive integer and $r\le v_2(M)$.
\end{corollary}
\begin{proof}
Exactness on constants requires $H(M)=0$, so $M$ must be an integer.
For an integer $M$, all derivatives of orders at most $r$ vanish at every
$Mn$, $n\ne0$, precisely when $r<1+v_2(M)$. The necessity of this
inequality is witnessed already by $n=1$ in \eqref{eq:alias}.
\end{proof}

\section{Classification of all exact signed phase filters}

\subsection{Multipliers of a finite signed measure}

Throughout this section a finite signed measure means a real, countably
additive Borel measure with finite total variation. For such a measure
$\nu$ on $\T$ of total mass one, set
\begin{equation}
 a_n=\int_\T e^{2\pi\ii n\varphi}\dd\nu(\varphi),\qquad n\in\Z.
 \label{eq:multipliers}
\end{equation}
Thus $a_0=1$, $a_{-n}=\overline{a_n}$, and
$|a_n|\le\norm{\nu}_{\TV}$. We say $\nu$ is \emph{uniformly exact through
$r$ at mesh $M$} when
\begin{equation}
 \Q^\nu_{M,\theta}(P)=\int_\R P(x)u(x)\dd x
 \quad\text{for all }P\in\Pol_r\text{ and all }\theta\in\T.
 \label{eq:uniform-exact}
\end{equation}
The same $\nu$ must work for all polynomials and all translations.

\begin{lemma}[Filtered Fourier test]
\label{lem:filtered-test}
Condition \eqref{eq:uniform-exact} is equivalent to
\begin{equation}
 a_n H^{(j)}(Mn)=0
 \quad (n\ne0,\ 0\le j\le r).
 \label{eq:multiplier-condition}
\end{equation}
Equivalently, $a_n$ may be nonzero only at indices for which $Mn$ is a
nonzero integer and $v_2(|Mn|)\ge r$.
\end{lemma}
\begin{proof}
Integrate the uniformly absolutely convergent series \eqref{eq:alias}
against $\nu$. Finite total variation justifies the interchange and gives
\begin{equation}
 \Def^\nu_{M,j}(\theta)
 :=\Q^\nu_{M,\theta}(x^j)-\mu_j
 =\left(\frac{\ii}{2\pi}\right)^j
 \sum_{n\ne0}a_nH^{(j)}(Mn)e^{2\pi\ii n\theta}.
 \label{eq:filtered-alias}
\end{equation}
A continuous periodic function vanishes identically if and only if all
its Fourier coefficients vanish. Exactness on the monomial basis is
therefore equivalent to \eqref{eq:multiplier-condition}. If $Mn$ is not an
integer, $H(Mn)\ne0$ already forces $a_n=0$. At a nonzero integer of
zero order $k$, all derivatives through $r$ vanish exactly when $k>r$;
otherwise the derivative of order $k$ is nonzero. Apply
Lemma~\ref{lem:kernel}.
\end{proof}

\subsection{Rational meshes: a cyclic symmetry is necessary and sufficient}

\begin{theorem}[Rational-mesh rigidity]
\label{thm:classification}
Let $M=a/b>0$, where $a,b\in\N_{>0}$ and $\gcd(a,b)=1$.
Let $r\ge0$ and put
\begin{equation}
 d=v_2(a),\qquad L=b\,2^{\max\{0,r-d\}}.
 \label{eq:L}
\end{equation}
For a finite signed Borel measure $\nu$ on $\T$ of mass one, the following
are equivalent:
\begin{enumerate}[label=(\roman*),leftmargin=*,itemsep=1pt]
\item $\nu$ is uniformly exact through degree $r$ at mesh $M$;
\item $a_n=0$ for every $n\notin L\Z$;
\item $\nu$ is invariant under rotation by $1/L$.
\end{enumerate}
For $L=1$, the rotation is the identity and (ii) imposes no restriction.
\end{theorem}

\begin{proof}
First identify the permitted frequencies in Lemma~\ref{lem:filtered-test}.
The condition $Mn\in\Z$ is equivalent to $b\mid n$ because $a/b$ is
reduced. Writing $n=bk$, the remaining condition is
\[
 v_2(|ak|)=d+v_2(|k|)\ge r.
\]
This says exactly that $2^{\max(0,r-d)}$ divides $k$, or $L$ divides $n$.
This proves (i)$\Leftrightarrow$(ii).

Let $R_h\nu$ be the pushforward of $\nu$ under $\varphi\mapsto\varphi+h$.
Its multiplier is $e^{2\pi\ii nh}a_n$. Invariance under $h=1/L$
therefore implies that $a_n=0$ whenever $L\nmid n$.
Conversely, if (ii) holds, $R_{1/L}\nu$ and $\nu$ have all the same
Fourier coefficients. Trigonometric polynomials are uniformly dense in
continuous functions on $\T$. Integrating their uniform approximations
against these finite-variation measures shows that the measures agree
on every continuous function, and hence are equal. This proves
(ii)$\Leftrightarrow$(iii).
\end{proof}

The theorem classifies singular continuous measures as well as atomic and
absolutely continuous ones. If $\nu$ has a density, that density must be
$1/L$-periodic almost everywhere. In the signed setting there is no
positivity premise hidden in the Fourier argument.

\begin{corollary}[Symmetrization description]
\label{cor:symmetrization}
At a rational mesh with the integer $L$ in \eqref{eq:L}, every measure
of the form
\begin{equation}
 \nu=\frac1L\sum_{j=0}^{L-1}R_{j/L}\rho,
 \qquad \rho(\T)=1,
 \label{eq:symmetrization}
\end{equation}
is uniformly exact through $r$. Conversely every exact measure has this
form, for example by taking $\rho=\nu$.
\end{corollary}
\begin{proof}
The averaging operator in \eqref{eq:symmetrization} has as its fixed
points precisely the $1/L$-rotation-invariant measures.
\end{proof}

\subsection{Irrational meshes: continuous averaging is unavoidable}

\begin{theorem}[Rational--irrational dichotomy]
\label{thm:irrational}
Let $M>0$ be irrational. A finite signed phase measure of mass one is
uniformly exact on constants if and only if it is normalized Haar measure
$m_\T$. Haar averaging, in fact, integrates $P(x)u(x)$ exactly for every
polynomial $P$. No finite atomic phase measure is uniformly exact on
constants at such a mesh.
\end{theorem}
\begin{proof}
For every nonzero integer $n$, $Mn$ is not an integer and so $H(Mn)\ne0$.
Lemma~\ref{lem:filtered-test}, even with $r=0$, gives $a_n=0$ for all
$n\ne0$. Together with $a_0=1$, these are the Fourier coefficients of
normalized Haar measure. Fourier uniqueness, as justified in the proof
of Theorem~\ref{thm:classification}, gives $\nu=m_\T$.

Conversely, directly unfolding the integral yields, for every integrable
function $f$,
\begin{equation}
 \int_0^1\frac1M\sum_{k\in\Z}
 f((k+\theta+\varphi)/M)\dd\varphi=\int_\R f(x)\dd x.
 \label{eq:haar-unfold}
\end{equation}
For a signed integrable function this follows by applying the positive
identity to its positive and negative parts. Take $f=Pu$. Haar measure
is nonatomic and cannot equal a finitely supported measure of mass one.
\end{proof}

\subsection{Minimum support, signed weights, and equality cases}

\begin{theorem}[Optimal atomic phase cost]
\label{thm:atomic}
Let $M=a/b$ and $L$ be as in Theorem~\ref{thm:classification}. Suppose
\begin{equation}
 \nu=\sum_{j=1}^K w_j\delta_{\varphi_j},\qquad
 w_j\in\R\setminus\{0\},\quad \sum_jw_j=1,
 \label{eq:atomic}
\end{equation}
where the phases are distinct modulo one. If $\nu$ is uniformly exact
through $r$, then $L$ divides $K$, and in particular $K\ge L$.
All such atomic measures are unions of complete $L$-cycles, with equal
weights within each cycle. When $K=L$ there is an $\alpha\in\T$ such that
\begin{equation}
 \nu=\frac1L\sum_{j=0}^{L-1}\delta_{\alpha+j/L}.
 \label{eq:optimal-filter}
\end{equation}
Thus every minimum-support exact signed filter is positive.
\end{theorem}
\begin{proof}
By Theorem~\ref{thm:classification}, rotation by $1/L$ leaves $\nu$
unchanged. The orbit of every point under this rotation has exactly
$L$ points. If an atom has nonzero mass $w$, all points in its orbit
are atoms of mass $w$. Distinct orbits are disjoint, so the support is
a disjoint union of $L$-point orbits and $L\mid K$.
If $K=L$, there is just one orbit, and its common mass is $1/L$
because the total mass is one. Conversely any union of such equally
weighted orbits with total mass one is invariant and hence exact.
\end{proof}

\begin{proposition}[The optimum is a finer lattice]
\label{prop:reindex}
The measure in \eqref{eq:optimal-filter} satisfies
\begin{equation}
 \Q^\nu_{M,\theta}(P)=\Q_{LM,\,L(\theta+\alpha)}(P).
 \label{eq:reindex}
\end{equation}
\end{proposition}
\begin{proof}
The left side is
\[
 \frac1{LM}\sum_{j=0}^{L-1}\sum_{k\in\Z}
 (Pu)\!\left(\frac{Lk+j+L(\theta+\alpha)}{LM}\right).
\]
Each integer has a unique representation $Lk+j$, $0\le j<L$.
Reindexing gives \eqref{eq:reindex}.
\end{proof}

For an integer mesh $M=2^dq$ with $q$ odd, a gain of $s\ge0$
translation-uniform degrees, from $d$ to $d+s$, costs exactly $2^s$
phases. More generally, with a budget of at most $K$ phases at the
rational mesh $a/b$, no uniform constant rule exists if $K<b$.
If $K\ge b$, the largest attainable uniform degree is
\begin{equation}
 r_{\max}=v_2(a)+\left\lfloor\log_2(K/b)\right\rfloor.
 \label{eq:budget}
\end{equation}
Indeed, the requirement $b2^{\max(0,r-d)}\le K$ is necessary by
Theorem~\ref{thm:atomic} and sufficient by \eqref{eq:optimal-filter}.

\begin{remark}[What this lower bound does not say]
The bound is on distinct \emph{phase shifts} of the fixed kernel rule,
with the weights fixed independently of the overall translation.
It is not a lower bound for arbitrary Gaussian quadrature nodes, nor for
a rule designed at just one phase, nor for weights allowed to vary with
$\theta$ or with the polynomial. These changes alter the problem's
quantifiers and invalidate the coefficientwise necessity argument.
\end{remark}

\section{A quantitative obstruction for under-budget signed filters}

The preceding support count concerns exact equality. A stronger statement
is possible: for fixed $M,r,K$ below the threshold, the entire class of
$K$-phase signed rules has a strictly positive lower bound on uniform
moment error. The weights need not lie in a compact or bounded set.

\subsection{A finite annihilator certificate}

\begin{lemma}[A nonzero early multiplier]
\label{lem:annihilator}
Let $\nu$ be as in \eqref{eq:atomic}. Then
\begin{equation}
 \max_{1\le n\le K}|a_n|\ge\frac1{2^K-1}.
 \label{eq:early-multiplier}
\end{equation}
More precisely, if $z_j=e^{2\pi\ii\varphi_j}$ and
\begin{equation}
 A(z)=\prod_{j=1}^K(z-z_j)=\sum_{n=0}^Kc_nz^n,
 \label{eq:annihilator-polynomial}
\end{equation}
then $|c_0|=1$, $\sum_{n=1}^K|c_n|\le2^K-1$, and
\begin{equation}
 \sum_{n=1}^Kc_na_n=-c_0.
 \label{eq:certificate}
\end{equation}
\end{lemma}
\begin{proof}
The polynomial $A$ vanishes at every $z_j$. Multiply $A(z_j)=0$ by
$w_j$ and sum. Since $a_0=1$, this gives \eqref{eq:certificate}.
The roots have modulus one, so $|c_0|=1$ and the elementary symmetric
coefficient bound is $|c_n|\le\binom Kn$. Hence
\[
 1\le\sum_{n=1}^K|c_n|\,|a_n|
 \le(2^K-1)\max_{1\le n\le K}|a_n|.
\]
No step uses positivity of the weights. In fact the conclusion also
holds for complex weights of total mass one.
\end{proof}

\subsection{Explicit uniform error bounds}

Define the dimensionless moment-error norm
\begin{equation}
 E_{M,r}(\nu)=\max_{0\le j\le r}\norm{\Def^\nu_{M,j}}_{L^\infty(\T)}.
 \label{eq:errornorm}
\end{equation}
The coordinate $x$ is already normalized to the fixed support $[-1,1]$.

\begin{theorem}[Quantitative under-budget obstruction]
\label{thm:quant-obstruction}
Let $M=a/b$ be rational, let $L$ be given by \eqref{eq:L}, and let
$1\le K<L$. For $1\le n\le K$ define
\begin{equation}
 j_n=\begin{cases}
 0,&Mn\notin\Z,\\
 1+v_2(Mn),&Mn\in\Z,
 \end{cases}
 \qquad
 B_n=(2\pi)^{-j_n}|H^{(j_n)}(Mn)|.
 \label{eq:Bn}
\end{equation}
Then $0\le j_n\le r$ and $B_n>0$. Every signed measure of
mass one supported on at most $K$ phases satisfies
\begin{equation}
 \boxed{\quad
 E_{M,r}(\nu)\ge
 \frac{\min_{1\le n\le K}B_n}{2^K-1}>0.
 \quad}
 \label{eq:quant-obstruction}
\end{equation}
For an irrational $M>0$ and every $K\ge1$, the corresponding conclusion
already holds on constants:
\begin{equation}
 \norm{\Def^\nu_{M,0}}_\infty
 \ge\frac{\min_{1\le n\le K}|H(Mn)|}{2^K-1}>0.
 \label{eq:irrational-quant}
\end{equation}
\end{theorem}

\begin{proof}
In the rational case, none of the indices $1,\ldots,K$ lies in $L\Z$.
The frequency calculation in Theorem~\ref{thm:classification} shows that
if $Mn$ is an integer, its zero order is at most $r$; otherwise order
zero is a nonvanishing witness. Thus \eqref{eq:Bn} is defined with
$j_n\le r$ and $B_n>0$.

If the support has fewer than $K$ points, multiply its annihilating
polynomial by additional factors with unit-modulus roots to reach degree
$K$. The coefficient proof of Lemma~\ref{lem:annihilator} remains valid.
The $n$th Fourier coefficient of the $j_n$th filtered defect has modulus
$B_n|a_n|$ by \eqref{eq:filtered-alias}. The modulus of a Fourier
coefficient of a continuous periodic function is at most its uniform
norm, so
\[
 E_{M,r}(\nu)\ge B_n|a_n|\qquad(1\le n\le K).
\]
Combine this with Lemma~\ref{lem:annihilator} to obtain
\eqref{eq:quant-obstruction}. More specifically, the actual nodes give
\begin{equation}
 E_{M,r}(\nu)\ge
 \left(\sum_{n=1}^K\frac{|c_n|}{B_n}\right)^{-1},
 \label{eq:node-dependent-bound}
\end{equation}
by inserting $|a_n|\le E_{M,r}(\nu)/B_n$ into
\eqref{eq:certificate}. Formula \eqref{eq:quant-obstruction} is a
node-independent consequence.

For irrational $M$, each $H(Mn)$ is nonzero. Use $j_n=0$ for every
$n$ and repeat the same argument. Positivity of the minimum follows
from taking a minimum over a nonempty finite set of strictly positive
numbers, not from a global lower bound for $H$.
\end{proof}

\begin{remark}[Explicit does not mean numerically large]
The constants can be very small, because $H$ decays rapidly and because
$2^K-1$ grows. The conclusion is a rigorous positive obstruction for
\emph{fixed} parameters, not a claim of a practical lower error floor
uniform in $M$ or $K$. Near a rational mesh, the irrational constant can
tend to zero. That behavior is consistent with the exact rational--
irrational dichotomy and prevents a false claim of a discontinuous
positive universal error threshold.
\end{remark}

\begin{corollary}[Integer-mesh certificate without high derivatives]
\label{cor:integer-B}
Suppose $M=2^dq$ with $q$ odd, $r=d+s$, and $K<2^s$.
For $1\le n\le K$, write $n=2^e\ell$ with $\ell$ odd and put
$j_n=d+e+1$. Then the constants in \eqref{eq:Bn} are
\begin{equation}
 B_n=\frac{j_n!\,|H(q\ell/2)|}{(2\pi Mn)^{j_n}}.
 \label{eq:explicit-B}
\end{equation}
\end{corollary}
\begin{proof}
At any positive integer $t=2^vq_0$ with $q_0$ odd, the leading-zero
calculation in Lemma~\ref{lem:leading-jet} below gives
$H^{(v+1)}(t)=-(v+1)!H(q_0/2)/t^{v+1}$.
Apply it to $t=Mn=2^{d+e}q\ell$.
\end{proof}

These certificates are finite in a useful sense: only the first $K$
phase multipliers and the corresponding nonzero kernel coefficients
are required. One need not search all Fourier frequencies to find a
witness that an under-budget rule fails.

\section{Sharper odd-dilation bounds and pointwise superconvergence}

\subsection{An inverse-square comparison}

The repository's multiple-angle argument isolates one exactly controlled
factor of the product and obtains an inverse first power
\cite{repoComb}. Two factors can be controlled exactly, uniformly in the
odd part of the mesh.

\begin{lemma}[Odd-dilation inverse-square bound]
\label{lem:dilation}
For all positive odd integers $q,n$,
\begin{equation}
 0<\left|\frac{H(qn/2)}{H(q/2)}\right|\le\frac1{n^2}.
 \label{eq:dilation}
\end{equation}
\end{lemma}
\begin{proof}
For every positive integer $n$ and real $x$,
$|\sin(nx)|\le n|\sin x|$. For example, this follows from
$e^{\ii nx}-e^{-\ii nx}=(e^{\ii x}-e^{-\ii x})
\sum_{j=0}^{n-1}e^{\ii(n-1-2j)x}$ and the triangle inequality.
Consequently $|\sinc(ny)|\le|\sinc(y)|$ for every real $y$.

Compare the factors of $H(qn/2)$ and $H(q/2)$.
At product index zero, $q$ and $qn$ are odd, so the absolute sine
numerators are both one. The factor ratio is exactly $1/n$.
At index one, the arguments of the sine are $\pi q/4$ and $\pi qn/4$;
for every odd integer $h$, $|\sin(\pi h/4)|=1/\sqrt2$.
This second factor ratio is also exactly $1/n$.
Every remaining factor ratio is at most one by the multiple-angle
inequality. No denominator factor vanishes. The comparison holds for
every finite product with at least two factors, and passing to the limit
gives the upper bound. Both full products are nonzero by
Lemma~\ref{lem:kernel}, which gives the strict lower inequality.
\end{proof}

The improvement is especially useful when a derivative of a trigonometric
harmonic costs one power of its frequency: the inverse square still leaves
a summable majorant at the lowest threshold degree. In fact the arithmetic
separation from zeros at every dyadic scale gives a stronger uniform result.

\begin{proposition}[A uniform multiscale odd-dilation envelope]
\label{prop:multiscale}
For positive odd $q$ and odd $n\ge3$, put $m=\lfloor\log_2 n\rfloor$.
Then
\begin{equation}
 \left|\frac{H(qn/2)}{H(q/2)}\right|
 \le \frac{2^{m(m+1)/2-1}}{n^{m+1}}
 \le \frac1{2\sqrt n}
       \exp\!\left(-\frac{(\log n)^2}{2\log2}\right).
 \label{eq:multiscale}
\end{equation}
In particular the relative odd-dilation ratio decreases faster than every
fixed inverse power of $n$, uniformly over all positive odd $q$.
\end{proposition}
\begin{proof}
At index $j\ge2$, oddness of $q$ gives
\[
 |\sin(\pi q/2^{j+1})|\ge\sin(\pi/2^{j+1})\ge2^{-j}.
\]
The first inequality follows by reducing the odd numerator modulo
$2^{j+1}$ and measuring its distance from the nearest multiple of that
integer. The second is $\sin x\ge2x/\pi$ on $[0,\pi/2]$.
Thus the absolute ratio of the $j$th factors is at most $2^j/n$,
as well as at most one by the multiple-angle inequality. Retain the
exact $n^{-2}$ from the first two factors and the bounds $2^j/n$
for $2\le j\le m$. All later ratios are at most one. Their product is
\[
 n^{-2}\prod_{j=2}^m\frac{2^j}{n}
 =2^{m(m+1)/2-1}n^{-m-1},
\]
with the product empty when $m=1$.
For the second inequality write $\log_2n=m+t$, $0\le t<1$.
The logarithm of this upper bound is
\[
 -\frac{(\log n)^2}{2\log2}-\frac12\log n-\log2
 +\frac{\log2}{2}(t^2-t),
\]
and the final term is nonpositive. The superalgebraic conclusion follows
because a negative quadratic in $\log n$ dominates every fixed negative
linear term in $\log n$.
\end{proof}

This is a relative estimate: it remains valid when $H(q/2)$ itself is
extremely small. An ordinary decay estimate for $H(t)$ at large $t$ would
not by itself provide this uniformity after dividing by $H(q/2)$. For the
root theorem below we retain the simpler inverse-square bound, which makes
the constants transparent; \eqref{eq:multiscale} can further reduce them.

\subsection{Leading jets at all integer zeros}

\begin{lemma}[The first nonzero derivative]
\label{lem:leading-jet}
Let $t_0=2^vq_0$ be a nonzero integer, $q_0$ odd, and put $k=v+1$.
Then
\begin{equation}
 H^{(k)}(t_0)=-\frac{k!}{t_0^k}H(q_0/2).
 \label{eq:leading-jet}
\end{equation}
In particular, if $M=2^dq$ with $q>0$ odd and $p=d+1$, then for
odd $n\in\Z$,
\begin{equation}
 H^{(p)}(Mn)=-\frac{p!}{(Mn)^p}H(qn/2),
 \label{eq:threshold-jet}
\end{equation}
whereas $H^{(p)}(Mn)=0$ for nonzero even $n$.
\end{lemma}
\begin{proof}
Split the product into its $k$ vanishing factors and the tail:
\[
 H(t)=\prod_{j=0}^{v}\sinc(t/2^j)\; H(t/2^{v+1}).
\]
At $t=t_0$, the linear coefficient of the $j$th factor is
$(-1)^{t_0/2^j}/t_0$. The product of their signs is
\[
 (-1)^{\sum_{j=0}^{v}t_0/2^j}
 =(-1)^{q_0(2^{v+1}-1)}=-1.
\]
Thus the coefficient of $(t-t_0)^k$ is
$-t_0^{-k}H(q_0/2)$, proving \eqref{eq:leading-jet}.
For even $n\ne0$, the zero order at $Mn$ is at least $d+2=p+1$,
so its $p$th derivative vanishes.
\end{proof}

This leading-jet identity is already represented in the repository's
\path{CombFirstDefect.lean} development \cite{repoComb}. Its derivation
is included to make the signs and constants of the new bounds auditable.

\subsection{A positive relative factor and complete zero sets}

Fix $M=2^dq$ with $d\ge0$ and $q>0$ odd. Write
\begin{equation}
 p=d+1,\qquad
 \tau_p(t)=\begin{cases}\sin t,&p\text{ odd},\\\cos t,&p\text{ even},\end{cases}
 \qquad \epsilon_p=(-1)^{\lfloor(p-1)/2\rfloor},
 \label{eq:tau}
\end{equation}
and define the nonzero amplitude
\begin{equation}
 A_{M,p}=\frac{2p!\,\epsilon_pH(q/2)}{(2\pi M)^p}.
 \label{eq:amplitude}
\end{equation}
Pairing positive and negative odd aliases in \eqref{eq:alias} gives
\begin{equation}
 \Def_{M,p}(\theta)=A_{M,p}
 \sum_{\substack{n\ge1\\n\ \mathrm{odd}}}
 \frac{H(qn/2)}{H(q/2)n^p}\,\tau_p(2\pi n\theta).
 \label{eq:paired}
\end{equation}
For instance, the sign factors for $p=1,2,3,4$ are respectively
$+,+,-,-$; this is a useful independent normalization check.

\begin{theorem}[Relative factorization and all integer-mesh phase zeros]
\label{thm:factorization}
With the notation above, there exists a continuous one-periodic real
function $G_{p,q}$ satisfying
\begin{align}
 \Def_{M,p}(\theta)
 &=A_{M,p}\tau_p(2\pi\theta)G_{p,q}(\theta),
 \label{eq:factorization}\\
 1-\eta_p\le G_{p,q}(\theta)&\le1+\eta_p,
 \qquad
 \eta_p=\sum_{\substack{n\ge3\\n\ \mathrm{odd}}}n^{-p-1}
 =(1-2^{-p-1})\zeta(p+1)-1.
 \label{eq:Gbound}
\end{align}
In particular $G_{p,q}$ is strictly positive, since
$\eta_p\le\pi^2/8-1<1$. The rule $\Q_{M,\theta}$ integrates
$\Pol_{d+1}$ exactly if and only if, modulo one,
\begin{equation}
 \theta\in Z_p:=
 \begin{cases}
 \{0,1/2\},&d\text{ even},\\
 \{1/4,3/4\},&d\text{ odd}.
 \end{cases}
 \label{eq:phase-set}
\end{equation}
At any other phase its degree of precision is exactly $d$.
\end{theorem}

\begin{proof}
For odd $n$ the quotients
\[
 R_{p,n}(t)=\frac{\tau_p(nt)}{\tau_p(t)}
\]
extend continuously across the zeros of the denominator. In the sine
case, the quotient is the polynomial $U_{n-1}(\cos t)$, where $U$ is
the Chebyshev polynomial of the second kind; equivalently the elementary
multiple-angle identity proves the same divisibility. In the cosine case,
$\cos(nt)$ is an odd polynomial in $\cos t$ because $n$ is odd,
so division by $\cos t$ also leaves a polynomial. The sine inequality
in the proof of Lemma~\ref{lem:dilation} gives $|R_{p,n}(t)|\le n$.
For cosines use the same inequality after shifting $t$ by $\pi/2$;
oddness of $n$ changes only a sign.

Define
\begin{equation}
 G_{p,q}(\theta)=1+
 \sum_{\substack{n\ge3\\n\ \mathrm{odd}}}
 \frac{H(qn/2)}{H(q/2)n^p}R_{p,n}(2\pi\theta).
 \label{eq:Gseries}
\end{equation}
By Lemma~\ref{lem:dilation}, the absolute value of each summand is at
most $n^{-p-1}$. Since $p\ge1$, the series is uniformly convergent.
It defines a continuous real periodic function and obeys
\eqref{eq:Gbound}. Multiplying it by $\tau_p(2\pi\theta)$ gives
\eqref{eq:paired}, including at the zeros by continuity. This proves
\eqref{eq:factorization}.

All degrees at most $d$ are exact at all phases by
Corollary~\ref{cor:unfiltered}. The only extra monomial to check for
$\Pol_{d+1}$ is $x^p$. Its defect vanishes exactly when
$\tau_p(2\pi\theta)=0$, because both $A_{M,p}$ and $G_{p,q}$ are
nonzero. These are exactly the phases in \eqref{eq:phase-set}.
Outside them, $x^{d+1}$ witnesses failure, proving the final assertion.
\end{proof}

\begin{remark}[The lowest-degree endpoint matters]
The theorem includes every odd mesh $M=q$, where $d=0$ and $p=1$.
A comparison of size $n^{-1}$ alone would bound the relative sine tail
by a divergent harmonic-type sum in this boundary case. The second
exact sinc ratio makes the relevant majorant $n^{-2}$ and settles the
boundary case by the same argument as every other integer mesh.
For higher degrees, the result also sharpens the quantitative constants
available from the weaker comparison. The statement does not assert
that the selected phases fail at the very next degree.
\end{remark}

\subsection{Simple zeros, phase conditioning, and arithmetic signs}

\begin{corollary}[Simple zeros with uniform relative derivative bounds]
\label{cor:simplezeros}
Every zero $\theta_*\in Z_p$ of $\Def_{M,p}$ is simple, and
\begin{equation}
 2\pi|A_{M,p}|(1-\eta_p)
 \le |\Def'_{M,p}(\theta_*)|
 \le2\pi|A_{M,p}|(1+\eta_p).
 \label{eq:root-slope}
\end{equation}
\end{corollary}
\begin{proof}
The defect is smooth by Lemma~\ref{lem:alias}. Divide
\eqref{eq:factorization} by $\theta-\theta_*$ and pass to the limit.
Continuity of $G_{p,q}$ suffices, and
$|\frac{\mathrm d}{\mathrm d\theta}\tau_p(2\pi\theta_*)|=2\pi$.
Now use \eqref{eq:Gbound}. No unproved differentiability of the quotient
series at a root is being assumed.
\end{proof}

\begin{corollary}[Global distance-to-phase bounds]
\label{cor:phase-distance}
Let $\rho=\operatorname{dist}_{\T}(\theta,Z_p)$, so $0\le\rho\le1/4$.
Then
\begin{equation}
 4|A_{M,p}|(1-\eta_p)\rho
 \le |\Def_{M,p}(\theta)|
 \le2\pi|A_{M,p}|(1+\eta_p)\rho.
 \label{eq:phase-distance}
\end{equation}
Consequently, an observed bound
$|\Def_{M,p}(\theta)|\le\varepsilon$ implies
\begin{equation}
 \operatorname{dist}_{\T}(\theta,Z_p)
 \le\frac{\varepsilon}{4|A_{M,p}|(1-\eta_p)}.
 \label{eq:phase-certificate}
\end{equation}
\end{corollary}
\begin{proof}
For the nearest zero, $|\tau_p(2\pi\theta)|=\sin(2\pi\rho)$.
On $[0,\pi/2]$, concavity gives $\sin x\ge2x/\pi$, while
$\sin x\le x$. Combine these with \eqref{eq:factorization}.
\end{proof}

The dependence on the amplitude is essential. The roots are well
conditioned after dividing by $A_{M,p}$, but the absolute defect can be
extremely small for large odd $q$. An unnormalized absolute error test
alone may then give a weak certificate for the actual phase.

\begin{proposition}[Binary sign rule]
\label{prop:sign}
Let $s_2(q)$ be the number of ones in the binary expansion of a positive
odd integer $q$. Then
\begin{equation}
 \sgn H(q/2)=(-1)^{s_2(q)-1}.
 \label{eq:binary-sign}
\end{equation}
For $\theta\notin Z_p$,
\begin{equation}
 \sgn\Def_{M,p}(\theta)
 =(-1)^{\lfloor(p-1)/2\rfloor+s_2(q)-1}
   \sgn\tau_p(2\pi\theta).
 \label{eq:defect-sign}
\end{equation}
\end{proposition}
\begin{proof}
The sign of the $j$th sinc factor at $q/2$ is
$(-1)^{\lfloor q/2^{j+1}\rfloor}$. All but finitely many are positive.
If $q=\sum_{h\ge0}b_h2^h$, then
\[
 \sum_{k\ge1}\lfloor q/2^k\rfloor
 =\sum_{h\ge1}b_h(2^h-1)=q-s_2(q).
\]
Since $q$ is odd, $q-s_2(q)$ has the same parity as $s_2(q)-1$.
This proves \eqref{eq:binary-sign}. Equation \eqref{eq:defect-sign}
follows from the positive factor in \eqref{eq:factorization}.
\end{proof}

The binary sign phenomenon itself is part of the existing repository
analysis \cite{repoComb}; the point here is its combination with the
sharper global magnitude and root estimates.

\section{Consequences for polynomial synthesis and concrete examples}

\subsection{Triangular moment deconvolution}

The repository transports quadrature identities into polynomial synthesis
using inverse moment operators \cite{repoLean}. The filter theorems respect
this transport exactly; they are not restricted to monomials as a preferred
application.

Define the operator on polynomials
\begin{equation}
 (TP)(y)=\int_\R P(y-x)u(x)\dd x
 =\sum_{j\ge0}\frac{(-1)^j\mu_j}{j!}P^{(j)}(y),
 \label{eq:moment-operator}
\end{equation}
where the sum is finite. Since $\mu_0=1$, $T$ preserves the leading
coefficient and is an invertible triangular operator on every $\Pol_r$.
Indeed $T=I+N$ with $N$ degree-lowering, so on $\Pol_r$,
$T^{-1}=\sum_{k=0}^r(-N)^k$.

\begin{proposition}[Exact transfer, including necessity]
\label{prop:transfer}
Fix $M>0$, a phase measure $\nu$, and $r\ge0$. The uniform quadrature
identity \eqref{eq:uniform-exact} is equivalent to
\begin{equation}
 P(y)=\int_\T\frac1M\sum_{k\in\Z}
 (T^{-1}P)\!\left(y-\frac{k+\theta+\varphi}{M}\right)
 u\!\left(\frac{k+\theta+\varphi}{M}\right)\dd\nu(\varphi)
 \label{eq:synthesis}
\end{equation}
for all $P\in\Pol_r$, all $y\in\R$, and all $\theta\in\T$.
\end{proposition}
\begin{proof}
If quadrature is exact, apply it to
$R_y(x)=(T^{-1}P)(y-x)$, whose degree is at most $r$.
Its integral against $u$ is $T(T^{-1}P)(y)=P(y)$.
Conversely fix any $y$. As $P$ ranges through $\Pol_r$, invertibility
of $T^{-1}$ and of the substitution $z\mapsto y-z$ implies that $R_y$
ranges through all of $\Pol_r$. Thus \eqref{eq:synthesis} is precisely
quadrature exactness for every such polynomial.
\end{proof}

Consequently the same cyclic invariance, minimum number of phases, and
irrational obstruction hold for this deconvolved synthesis scheme.
The selected-phase criterion in Theorem~\ref{thm:factorization} likewise
transfers to a single phase at the first threshold degree.

\subsection{Phase-cost examples}

\begin{example}[An odd denominator]
Take $M=3/5$. Then $v_2(3)=0$.
Uniform exactness through degree two requires
$L=5\cdot2^2=20$ phases, even with signed weights. At the minimum,
the phases are $\alpha+j/20$ with weights $1/20$, and the resulting
rule is the mesh-$12$ rule with phase $20(\theta+\alpha)$.
Since $v_2(12)=2$, the degree matches exactly. A rule with at most
19 phases has a positive uniform moment-error lower bound supplied
by Theorem~\ref{thm:quant-obstruction}.
\end{example}

\begin{example}[Integer mesh versus special phase]
For $M=12$, $d=2$ and $p=3$.
Every phase integrates quadratics. Exactly phases $0$ and $1/2$
modulo one integrate cubics as well. To integrate cubics for every
overall translation using a fixed filter requires two antipodal
phases of equal weight. Uniform degree five requires eight phases;
at the minimum this is the mesh-$96$ rule. The phase count is about
translation-uniform exactness, not about the two exceptional cubic phases.
\end{example}

\begin{example}[A concrete positive obstruction]
Take $M=1$, target degree $r=1$, and $K=1$. Here $L=2$, $j_1=1$, and
\[
 E_{1,1}(\nu)\ge\frac{|H(1/2)|}{2\pi}.
\]
Since a one-point signed measure of mass one has weight one, this is a
lower bound valid for every translation of that single phase. At a
\emph{fixed} phase, however, the linear moment is exact at $0$ and $1/2$. This example explicitly
separates the quantifiers responsible for the obstruction.
\end{example}

\begin{example}[Irrational scale]
For $M=\sqrt2$ and any $K$-phase signed filter,
\[
 \sup_\theta|\Q^\nu_{\sqrt2,\theta}(1)-1|
 \ge\frac{\min_{1\le n\le K}|H(n\sqrt2)|}{2^K-1}>0.
\]
No finite filter can produce a partition of unity at all translations.
Continuous uniform phase averaging does so exactly by
\eqref{eq:haar-unfold}.
\end{example}

\section{Exact evaluation and reproducible verification}

\subsection{A moment recurrence}

The verification program does not use a truncated infinite product to
certify rational quadrature values. It evaluates dyadic samples exactly
using finite convolution cells and rational moments.

\begin{lemma}[Rational moments]
\label{lem:moments}
The moments satisfy $\mu_0=1$, $\mu_{2n+1}=0$, and for $n\ge1$,
\begin{equation}
 \mu_{2n}=\frac1{2^{2n}-1}
 \sum_{j=1}^n\binom{2n}{2j}\frac{\mu_{2n-2j}}{2j+1}.
 \label{eq:moments}
\end{equation}
In particular all moments are rational, with
\[
 \mu_2=\frac19,\qquad \mu_4=\frac{19}{675},\qquad
 \mu_6=\frac{583}{59535}.
\]
\end{lemma}
\begin{proof}
The random series has the distributional identity
$X\overset{\mathrm d}=(V+X')/2$, where $V$ is uniform on $[-1,1]$
and independent of the copy $X'$ of $X$. Expand the $2n$th power.
Odd moments vanish by symmetry and $\mathbb E V^{2j}=1/(2j+1)$.
Moving the term involving $\mu_{2n}$ to the left gives
\eqref{eq:moments}.
\end{proof}

\subsection{A finite exact dyadic evaluator}

For a nonnegative integer $m$, write $\sigma(m)=(-1)^{s_2(m)}$.
For $N\ge1$ and real $x$, put
\begin{equation}
 a_m(x)=x+1-2^{-N}-m2^{-(N-1)},\qquad 0\le m<2^N.
 \label{eq:cell-offset}
\end{equation}

\begin{proposition}[Exact dyadic evaluation formula]
\label{prop:dyadic-evaluator}
Suppose $x$ is dyadic with $|x|<1$, and write its reduced denominator
as $2^d$. Set $N=d+1$. Then
\begin{equation}
 \begin{split}
 u(x)=\frac{2^{N(N-1)/2}}{(N-1)!}
 \sum_{\substack{0\le m<2^N\\a_m(x)>0}}\sigma(m)
 \sum_{\substack{0\le j\le N-1\\j\ \mathrm{even}}}
 \binom{N-1}{j}
 a_m(x)^{N-1-j}2^{-Nj}\mu_j.
 \end{split}
 \label{eq:exact-evaluator}
\end{equation}
The value is therefore obtainable with finite rational arithmetic.
\end{proposition}

\begin{proof}
Let $u_N$ be the density of $\sum_{j=0}^{N-1}2^{-j}U_j$.
Convolution of the $N$ uniform densities gives the finite
inclusion--exclusion formula
\begin{equation}
 u_N(t)=\frac{2^{N(N-1)/2}}{(N-1)!}
 \sum_{m=0}^{2^N-1}\sigma(m)
 \bigl(t+1-2^{-N}-m2^{-(N-1)}\bigr)_+^{N-1}.
 \label{eq:box-spline}
\end{equation}
To verify this, shift the uniforms to intervals $[0,2^{-j}]$.
The density of their sum is the derivative of the box-truncated simplex
volume: inclusion--exclusion subtracts subsets of upper-bound violations,
with sign $(-1)^{|S|}$ and displacement $\sum_{j\in S}2^{-j}$.
These subset sums are exactly $m2^{-(N-1)}$, and $|S|$ has parity
$s_2(m)$. The reciprocal product of interval lengths is
$2^{N(N-1)/2}$. This yields \eqref{eq:box-spline}.

The remaining random tail is $2^{-N}X'$, independent of the first
$N$ terms. Hence $u(x)=\mathbb E u_N(x-2^{-N}X')$.
Because $2^Nx$ is even, each $2^Na_m(x)$ is an odd integer. Thus
$|a_m(x)|\ge2^{-N}$. As $|X'|\le1$, the argument
$a_m(x)-2^{-N}X'$ never changes sign except possibly at an endpoint
of probability zero. The negative terms vanish; the positive terms
may be expanded as ordinary powers over the whole expectation.
Only the even moments survive, and the binomial expansion is precisely
\eqref{eq:exact-evaluator}. At $N=1$ the truncated zeroth power is
interpreted as the positive-half-line indicator; the same almost-sure
argument applies.
\end{proof}

The algorithm is a transparent exact certificate, not a claim of optimal
bit complexity. Its number of summands grows exponentially with the
denominator depth. More efficient evaluation is a separate algorithmic
question and is not needed for the proofs above.

\subsection{Reproduced exact values and scope of testing}

The bundled \path{verify.py} uses Python's arbitrary-precision integers
and \path{fractions.Fraction} for the exact part. The default run checks
all eight phases $h/8$, $0\le h<8$, for meshes $2^d$, $0\le d\le4$;
all lower-degree moments; the precise first-threshold zero pattern;
and cyclic-filter reindexing examples. The arithmetic criterion for
permitted frequencies is checked for reduced $a/b$ with
$1\le a\le16$, $1\le b\le12$, $0\le r\le6$, and $1\le n\le64$.

\begin{center}
\begin{tabular}{rrrr}
\toprule
$M$ & $p=v_2(M)+1$ & $\Def_{M,p}(0)$ & $\Def_{M,p}(1/4)$\\
\midrule
1&1&0&$13/72$\\[2pt]
2&2&$1/72$&0\\[2pt]
4&3&0&$-4643/11059200$\\[2pt]
8&4&$-23/5529600$&0\\[2pt]
16&5&0&$4966069/383551316951040$\\
\bottomrule
\end{tabular}
\end{center}
The four even-mesh rows reproduce values already in the repository's
comb manuscript \cite{repoComb}. They are cross-checks, not new numerical
claims. The exact run passes 261 quadrature, moment, and reindexing
assertions and 53,760 frequency-arithmetic assertions.

An optional \path{mpmath} part, run at 75 decimal digits, checks 240
odd-dilation ratios and their multiscale envelopes, 36 leading-zero
derivative formulas, and 12
Fourier-normalization comparisons with the independent exact evaluator.
The largest absolute discrepancy in the last group, using 200 positive
odd harmonics, is approximately $6.40\times10^{-19}$. These are
floating-point diagnostics, \emph{not} interval enclosures. The numerical
criterion used in the program is $10^{-18}$. No universal mathematical
conclusion in this article depends on passing those finite tests.

For orientation, the first relative-tail constants are
\begin{center}
\begin{tabular}{rll}
\toprule
$p$ & $\eta_p$ (approximately) & Guaranteed range of $G_{p,q}$\\
\midrule
1&0.2337005501&$[0.7662994499,\,1.2337005501]$\\
2&0.0517997903&$[0.9482002097,\,1.0517997903]$\\
3&0.0146780316&$[0.9853219684,\,1.0146780316]$\\
4&0.0045237628&$[0.9954762372,\,1.0045237628]$\\
\bottomrule
\end{tabular}
\end{center}
The exact guaranteed endpoints are $1\pm\eta_p$; the displayed decimal
endpoints are rounded illustrations, not outward-rounded certificates.

\section{Formalization map and proof obligations}

The article is independent of any unverified repository theorem because
its kernel facts, Fourier criterion, and all new conclusions have explicit
proofs. Nevertheless the existing Lean modules provide a natural route
to formalizing the extensions. This section is a plan, not a report of a
successful Lean build.

The inspected \path{RvachevSuperconvergentSynthesis.lean} imports
\path{CombDefectSeries} and \path{RvachevPolynomialSynthesis} and exposes
physical-coordinate quadrature and polynomial-deconvolution statements
\cite{repoLean}. The canonical prose also names existing formal
multiple-angle and first-defect lemmas \cite{repoComb}. These are useful
interfaces to reuse after checking their current types at the pinned
revision.

\begin{center}
\begin{tabular}{p{0.28\linewidth}p{0.64\linewidth}}
\toprule
Proposed layer & Main formal obligations\\
\midrule
Inverse-square comparison & Isolate the two factors at $q/2$ and $q/4$;
prove their absolute numerator identities for odd integers; pass the finite
comparison to the nonzero infinite product.\\[4pt]
Relative defect factor & Define removable sine/cosine quotients for odd
frequencies; establish the $n$ bound; use a uniform summability theorem
and transfer the exact alias identity.\\[4pt]
Signed phase filters & Represent finite signed measures on the additive
circle; justify termwise integration against finite total variation;
prove Fourier uniqueness through trigonometric density.\\[4pt]
Arithmetic rigidity & Formalize the reduced-fraction condition
$b\mid n$ and the subsequent $2$-adic divisibility threshold.\\[4pt]
Atomic optimality & Show that invariance permutes nonzero atoms in
orbits of length $L$ and preserves the mass of each atom.\\[4pt]
Quantitative obstruction & Form the annihilating polynomial on unit-modulus
phase nodes, bound its elementary-symmetric coefficients, and compare
Fourier coefficients with uniform norms.\\
\bottomrule
\end{tabular}
\end{center}

Several distinctions should remain visible in a formal statement.
First, $M$ is a positive real in the analytic definition and only later
specialized to an integer or reduced rational. Second, $r$ may be zero;
this is essential to the irrational obstruction. Third, phase equality
is modulo one, rather than equality with just one of four literal real
representatives. Fourth, support size counts distinct nonzero atoms;
zero weights and duplicate listed phases must be removed before the
sharp cardinality assertion. Finally, exactness is simultaneously for
all $P\in\Pol_r$ and all $\theta$. Dropping the latter quantifier
produces a different theorem and cannot be hidden in an abbreviation.

\section{Further research questions and topics}

The following questions are not asserted to be globally open in every
formulation. They are precise next tasks suggested by the proved results,
and should receive their own literature checks before a novelty claim.
The first is closely related to a mixed-jet nonvanishing issue retained
in the repository's comb discussion \cite{repoComb}.

\begin{question}[The next defect at a selected phase]
For $M=2^dq$ and $\theta\in Z_{d+1}$, determine exactly when
$\Def_{M,d+2}(\theta)=0$. Is the degree of precision always exactly
$d+1$, or do additional exceptional meshes and phases occur?
\end{question}
The first-defect theorem does not answer this: the next derivative
involves higher jets at the odd aliases and leading jets at an additional
valuation layer. A proof of noncancellation has to control both. For
example, at $M=1$ the two selected phases fail on $x^2$, with defects
$-1/9$ and $5/36$, respectively, but finite examples do not settle the
general question.

\begin{question}[Sharp relative dominance and root conditioning]
Determine the best constant $\eta_p^*$ such that
$|G_{p,q}(\theta)-1|\le\eta_p^*$ for all positive odd $q$ and all
$\theta$. Is the supremum attained, and how does it depend on $p$?
\end{question}
Our $\eta_p$ comes from a termwise triangle inequality. It does not exploit
the arithmetic correlation of the signs and amplitudes across distinct
odd harmonics. Sharp constants could improve certified phase recovery.

\begin{question}[Optimal quantitative under-budget error]
For a rational mesh and $K<L$, compute or estimate sharply
\[
 \inf\{E_{M,r}(\nu):\nu\text{ has mass one and at most }K
 \text{ atoms}\}.
\]
How do positive and signed filters differ away from exactness?
\end{question}
Theorem~\ref{thm:quant-obstruction} proves strict positivity without
compactness assumptions on signed weights. It does not identify a
minimizer or a sharp constant. Coalescing phases and unbounded signed
weights are genuine issues in an optimization proof, even though they
cannot drive the error to zero below the threshold.

\begin{question}[Arithmetic behavior near rational scales]
For fixed $K$, describe the smallness of the irrational lower bound
in \eqref{eq:irrational-quant} as $M$ approaches a rational mesh. Can
its vanishing rate be expressed using the multiplicities of the
nearby integer zeros and rational approximation of $M$?
\end{question}
The leading-jet identity suggests a local power law near any fixed
resonance. A uniform theorem over families with growing $K$ would also
need control of which resonance supplies the minimum.

% ed. (2026-09-30): the fixed-K question is answered by a later draft of this
% directory; uniformity in K and the constants stay open.
\begin{ednote}
For fixed $K$ this question is answered by Theorem \texttt{thm:resonance}
of the later unreviewed draft \path{Approximate_Phase_Rigidity/} of this
directory (2026-09-30), for the optimal error itself and not only for the
bound \eqref{eq:irrational-quant}. Near a reduced rational $a/b$, the
infimum of $E_{M,r}(\nu)$ over mass-one filters with at most $K$ atoms,
positive, real signed or complex, with nodes and weights allowed to depend
on $M$, is of exact order
$|M-a/b|^{\max\{0,\,v_2(a)+1+\lfloor\log_2(K/b)\rfloor-r\}}$ for $K\ge b$,
and stays bounded away from zero for $K<b$. The exponent is the largest zero
multiplicity of $H$ at the frequencies $na/b$, $1\le n\le K$, minus $r$; for
$r=0$ the same multiplicity gives the vanishing order of the minimum in
\eqref{eq:irrational-quant}. An explicit regular polygon attains the order.
The constants and the neighbourhood depend on $a/b$, $K$ and $r$; their
uniformity as $K$ grows, and the optimal leading constants, remain open
there. That draft does not compute the under-budget infimum of the preceding
question at the rational mesh itself.
\end{ednote}

\begin{question}[Several dimensions and correlated phases]
For a product of Rvachev kernels, classify phase measures giving uniform
exactness on rectangular or total-degree polynomial spaces. Can correlated
shifts outperform the Cartesian product of one-dimensional optimal rules?
\end{question}
One must not simply tensor the one-dimensional necessity statement.
A multidimensional alias is a product of derivative factors; a zero
factor in one coordinate can annihilate the whole alias. The permitted
Fourier set therefore has a different combinatorial structure, and the
one-dimensional cyclic-invariance proof alone does not imply a product
support lower bound. This distinction makes the correlated problem a
substantive next target.

\begin{question}[Other dilation systems]
Replace the dyadic weights in the random series by $b^{-j}$ for an
integer $b>2$, or by a nonintegral geometric ratio. Which phase
classification, exact-filter symmetry, and finite-support rigidity
statements survive?
\end{question}
The measure-level Fourier test survives unchanged for a smooth compactly
supported kernel. The hard part is the arithmetic geometry of its zero
set and multiplicities. The particularly strong two-factor half-integer
comparison used here need not persist in another base.

\begin{question}[Finite certificates and efficient formalization]
Construct a certified implementation that returns either an exact cyclic
symmetry certificate or a nonzero Fourier witness for a proposed filter,
with rational or rigorously enclosed transcendental phase data. Formalize
the main signed-filter theorem and the under-budget bound in Lean.
\end{question}
The finite annihilator supplies a natural witness search bound: the first
$K$ multipliers suffice below the exactness threshold. The dyadic exact
evaluator supplies a separate rational route for regression tests.
Neither ingredient alone proves a favorable end-to-end bit-complexity
bound; that would require a specified representation of phases and
controlled evaluation of the nonzero kernel coefficients.

\section{Conclusion}

Uniform phase averaging has an exact arithmetic structure. At a reduced
rational mesh $a/b$, polynomial degree $r$ forces invariance under a
cyclic group of order $b2^{\max(0,r-v_2(a))}$. At the minimum support
size, the only signed filter is the positive equally spaced one, and
its action is exactly ordinary mesh refinement. At an irrational mesh,
only continuous Haar averaging can be exact even on constants for all
translations. A finite annihilating polynomial turns these structural
facts into positive lower error bounds for every under-budget signed
rule.

The complementary inverse-square odd-dilation estimate and its uniform
log-Gaussian strengthening refine the repository's established
first-harmonic approach. It gives a single
integer-mesh theorem with explicit relative constants, simple phase
zeros, and a global distance-to-zero certificate, without confusing
known phase classifications with new results or claiming unproved
higher-order maximality. The principal remaining opportunities are
sharp approximation below the phase threshold, mixed higher jets,
and genuinely multidimensional phase correlations.

\begin{thebibliography}{9}

\bibitem{arias}
J.~Arias de Reyna,
\emph{An infinitely differentiable function with compact support:
Definition and properties}, arXiv:1702.05442, 2017.
English translation of the author's 1982 article.
\url{https://arxiv.org/abs/1702.05442}.

\bibitem{zakharov}
V.~G.~Zakharov,
\emph{Polynomial spaces reproduced by elliptic scaling functions},
arXiv:1310.8425, 2013.
\url{https://arxiv.org/abs/1310.8425}.

\bibitem{repoLean}
V.~Reshetnikov and repository contributors,
\emph{ProveIt: RvachevSuperconvergentSynthesis.lean},
\path{Analysis/FabiusFunction/Lean/FabiusFunction/RvachevSuperconvergentSynthesis.lean},
pinned commit \snapshot, inspected September 28, 2026.
\url{https://github.com/VladimirReshetnikov/ProveIt/blob/37e61c1fdec28c6e7ab7ff445043993077b42cbd/Analysis/FabiusFunction/Lean/FabiusFunction/RvachevSuperconvergentSynthesis.lean}.

\bibitem{repoComb}
V.~Reshetnikov and repository contributors,
\emph{Comb Interpolation and Sampling Frontiers}, chapter
\path{03_additive_dyadic.tex}, pinned commit \snapshot.
Relevant source ranges inspected: lines 1--100, 1500--1650,
2030--2200, and 2200--2420. In particular:
\path{lem:multiple-angle}, \path{thm:phase-zero-set},
\path{thm:up-composite-mesh}, and the discussion of mixed-jet nonvanishing.
\url{https://github.com/VladimirReshetnikov/ProveIt/blob/37e61c1fdec28c6e7ab7ff445043993077b42cbd/Analysis/FabiusFunction/docs/semi-formalized-research-frontiers/drafts/inverse-and-sampling/comb-interpolation/comb_interpolation_synthesis/chapters/03_additive_dyadic.tex}.

\end{thebibliography}

\end{document}
